% 写作模板，不是已完成的课程报告。
% 用本人核验过的真实工作替换所有 \todo{...}，并删除教学提示。
% 全文、参考文献及声明合计不得超过五页；不得添加独立封面或附录。
% 标题、成员信息、参考文献及 AI/贡献声明由主模板集中生成。

\section{问题与动机}
\todo{说明野外植物病害识别的问题、实际意义，以及九种宿主子集的研究范围。
明确所选 Option，解释成功结果能够说明什么。结论应与实际证据相符。}

\ifnum\ReportOption=2
\todo{在训练前提出可证伪的假设：二值病斑监督如何可能改善未见过的
宿主与病害组合的拒识，同时保持已知类别识别能力。预先定义并论证已知类别
macro-F1 的可接受下降幅度，说明何种证据支持、反驳或不足以判定该假设。}
\else
\todo{定义闭集分类任务：输入单张 RGB 照片，输出宿主与病害组合类别。
说明症状相似、病斑较小、背景变化及类别不平衡带来的挑战。}
\fi

\section{相关工作}
课程指定数据版本为 PlantSeg v7 \citep{wei2025plantsegv7}，相关数据描述论文
提供数据集背景 \citep{wei2026plantseg}。
\todo{把实际阅读过的方法加入 references.bib，说明其训练数据、监督方式、
推理输入和评估协议与本研究的关系。核对每条文献信息，并为借鉴的实质性思路提供引用。}

\ifnum\ReportOption=2
\todo{讨论至少三项密切相关的方法，包括未知植物病害识别以及病斑引导或
grounded recognition。记录当前文献搜索的日期和范围；区分拒绝未知类别、
预测未见类别的身份和对健康图像拒答。说明本研究是复现、扩展还是新的比较，
并为任何新颖性主张提供证据。}
\fi

\section{数据与实验协议}
\todo{仅使用 2025 年 11 月 25 日发布的 PlantSeg v7，保留 apple、banana、
citrus、corn、potato、soybean、tomato、wheat 和 zucchini。
从下载的元数据中推导实际类别标识与样本数，把每个宿主与病害组合视为一类。
保留原始标签，以及约 70/10/20 的训练、验证、测试划分。}

\todo{引用冻结的数据清单，报告各分区的类别数、无法读取的文件、重复检查与排除记录。
说明如何处理相同图像及可能的近重复泄漏。仅用训练数据拟合数据相关预处理，
增强图像随原训练图像保留在同一分区。给出固定类别顺序、未定义指标的处理方式，
以及没有评估样本的类别如何参与统计。}

\ifnum\ReportOption=2
\todo{记录两套冻结的类别角色映射。所有 citrus 类为开发未知类，所有 potato 类为测试未知类。
其余每种宿主的病害标签按字母排序：Assignment A 将第一、第二类分别设为开发、测试未知类；
Assignment B 则使用第二、第三类。仅在已知类上训练，丢弃保留类别中不允许使用的图像。
分别报告熟悉宿主的未知组合与陌生宿主对照；样本量允许时，区分共享病害标签与未见病害标签。}
\fi

\section{方法}
\todo{描述网络结构、初始化来源或精确的预训练权重标识、RGB 预处理、增强、损失、
优化器、学习率、训练预算及停止规则，并论证缩放与裁剪不会系统性移除病害症状。
记录随机种子、依赖版本、硬件、训练耗时及 checkpoint 选择规则。
验证与测试使用确定性变换、评估模式和关闭梯度。
不得使用在 PlantSeg 或重叠 PlantWild 数据上训练的 checkpoint，也不得把含类别的文件名
或元数据作为预测输入。}

\ifnum\ReportOption=2
\todo{解释病斑监督的具体机制，以及它为何可能有效或失效。
仅在训练阶段使用已知训练类别的二值病斑掩码；所有评估预测仅输入 RGB。
构造不使用病斑监督的匹配分类器，控制骨干网络、初始化、预处理、增强以及训练和调参预算。
对声称有贡献的附加组件分别进行移除消融。}
\else
\todo{描述多数类预测器和基础 CNN，设计至少两项受控改进，并尽量一次只改变一个因素。
说明如何根据验证集 macro-F1 选择最终配置，以及如何在最终测试前冻结所有选择。}
\fi

\section{实验与结果}
\todo{使用 notebook 生成并核验的真实结果替换表~\ref{tab:results} 中的全部横线，
按要求报告样本量与不确定性。横线仅表示占位，不能解释为零或未报告的已完成实验。
一致地使用 macro-F1（宏平均 F1）等指标定义。}

\begin{table}[htbp]
\centering
\small
\ifnum\ReportOption=2
\caption{单套类别映射的结果模板。所有数值均为占位，不代表已进行的实验或实测结果。}
\label{tab:results}
\begin{tabularx}{\linewidth}{@{}lXXX@{}}
\toprule
模型 & 熟悉宿主 AUOSCR & 已知类 macro-F1 & 未知类 AUROC \\
\midrule
匹配基线 & -- & -- & -- \\
病斑监督 & -- & -- & -- \\
组件消融 & -- & -- & -- \\
\bottomrule
\end{tabularx}
\else
\caption{测试结果模板。所有数值均为占位，不代表已进行的实验或实测结果。}
\label{tab:results}
\begin{tabularx}{\linewidth}{@{}lXXX@{}}
\toprule
系统 & Macro-F1 & Top-1 accuracy & Balanced accuracy \\
\midrule
多数类预测器 & -- & -- & -- \\
基础 CNN & -- & -- & -- \\
受控改进 1 & -- & -- & -- \\
受控改进 2 & -- & -- & -- \\
\bottomrule
\end{tabularx}
\fi
\end{table}

\ifnum\ReportOption=2
\todo{两套类别映射各运行三个训练种子，即每个模型配置六次训练；各映射独立训练。
按映射分别报告种子均值、标准差及方法间配对差值。
比较 maximum-softmax、maximum-logit 和 energy 分数；仅用基线验证结果选择主要拒识规则，
并将同一规则用于两模型的主比较。报告熟悉宿主 AUOSCR、已知类 macro-F1/accuracy、
未知类 AUROC，并独立呈现陌生宿主结果。}
\todo{根据已知类验证 macro-F1 选择 checkpoint；在预设容差下，根据熟悉宿主验证 AUOSCR
选择配置。在查看任一映射的测试结果之前冻结两套映射的所有选择。
仅用已知验证图像确定约 95\% 接受率的阈值，再报告测试已知接受率、正确且被接受的比例，
以及未知类误接受率。}
\else
\todo{报告基线与改进系统的最终测试指标，以及逐类 precision、recall、F1、support 和混淆矩阵。
若报告仅显示可读的类别子集，notebook 中仍需保留完整矩阵。
结合实际证据解释各项受控改进的作用。}
\fi

\begin{figure}[htbp]
\centering
\fbox{\parbox[c][23mm][c]{0.9\linewidth}{\centering
图示占位 / Figure placeholder\\
此处没有数据，也不表示任何实测结果。\\
请替换为 notebook 生成的有效图示。}}
\caption{可视化占位，并非实验结果。请将方框及本图注替换为提交 notebook 生成、
带编号和清晰标签的真实图示。}
\label{fig:diagnostic}
\end{figure}

\todo{将图~\ref{fig:diagnostic} 替换为有解释价值的图，例如混淆热图、标记选中 checkpoint
的学习曲线、标注真实类别与预测的失败案例，或 Option 2 的 OSCR 曲线。
提供清楚的坐标轴、颜色标尺、类别标签及相关样本数。
提交的 notebook 必须能重新生成报告中的每一张数据派生图和每一张结果表。}

\section{解释与局限性}
\todo{解释代表性的成功与失败案例、类别不平衡、潜在泄漏、数据和适用范围限制，以及计算约束。
区分观测结果与推测，不将结论推广至未经评估的群体。说明证据支持的结论，
并提出一项有理由的后续研究。}
\ifnum\ReportOption=2
\todo{结合 AUOSCR 与已知类 macro-F1 的配对差值、种子与映射间变异、各类未知组的表现，
评估训练前的假设。混合或负面结果也是有效结论。
仅包含 potato 的陌生宿主对照不能支持对所有陌生植物宿主的泛化主张。}
\fi

% 参考文献及 AI/团队贡献声明由主模板集中渲染，勿在此重复。
